\documentclass[twocolumn,10pt]{article}
\usepackage[utf8]{inputenc}
\usepackage{geometry}
\geometry{letterpaper, margin=0.75in}
\usepackage{cite}
\usepackage{hyperref}
\usepackage{graphicx}
\usepackage{listings}
\usepackage{xcolor}

\title{A Layered Security Architecture for the FIX Protocol: \\ Integrating TLS 1.3 and AES-256-GCM}
\author{Gemini CLI Stack Implementation Report \\ College of Engineering and Technology}
\date{April 14, 2026}

\begin{document}

\maketitle

\begin{abstract}
The Financial Information eXchange (FIX) protocol is the industry standard for real-time electronic communication in financial markets. While natively unencrypted, modern trading environments require robust security. This paper describes a modular C++ implementation of a FIX stack that employs a dual-layer security approach: transport-level encryption via TLS 1.3 and application-level payload encryption using AES-256-GCM. We demonstrate a functional round-trip between a simulator and an exchange, ensuring data integrity and confidentiality.
\end{abstract}

\section{Introduction}
The FIX protocol \cite{fix42} defines a series of messaging semantics for trade-related communication. Historically, FIX sessions relied on private leased lines for security. However, with the shift toward cloud-based trading and public internet transit, application-layer and transport-layer security have become mandatory. This project implements a minimal but complete FIX stack (version 4.2) in C++ to explore these security paradigms.

\section{System Architecture}
The system is divided into five distinct layers to ensure separation of concerns:
\begin{enumerate}
    \item \textbf{FIX Engine:} Manages session state, sequence numbers, and heartbeats.
    \item \textbf{Codec:} Performs serialization and parsing of FIX tag-value pairs.
    \item \textbf{Security Layer:} Handles optional AES-256-GCM payload encryption.
    \item \textbf{Transport:} Manages TCP sockets wrapped in TLS 1.3.
    \item \textbf{Exchange Simulator:} A basic matching logic to process orders.
\end{enumerate}

\section{Security Implementation}

\subsection{Transport Layer Security (TLS)}
The stack utilizes TLS 1.3 \cite{rfc8446} to provide a secure tunnel for all communication. This is implemented using the OpenSSL library \cite{openssl}. The server employs a self-signed RSA-4096 certificate. TLS ensures that the entire FIX stream is invisible to network sniffers and protects against man-in-the-middle attacks.

\subsection{AES-256-GCM Payload Encryption}
In addition to TLS, we implemented an application-level encryption layer using AES in Galois/Counter Mode (GCM) \cite{nist80038d}. GCM was selected because it provides both confidentiality and data authentication (AEAD). 

The encryption framing prepends a 12-byte random nonce to the ciphertext. The FIX body (excluding tags 8, 9, and 10) is encrypted, and the \texttt{BodyLength} (tag 9) and \texttt{CheckSum} (tag 10) are updated to reflect the encrypted payload size and content. This ensures that even if TLS is terminated at a proxy, the sensitive order data remain encrypted until they reach the final engine.

\section{Validation and Results}
The implementation was verified using a Python-based simulator (\texttt{simplefix}). A full round-trip flow was observed:
\begin{enumerate}
    \item \textbf{Logon (35=A):} Successful handshake and TLS termination.
    \item \textbf{Order (35=D):} Execution of a \texttt{NewOrderSingle} with randomized symbols.
    \item \textbf{Report (35=8):} Reception of an \texttt{ExecutionReport} indicating a filled order.
    \item \textbf{Logout (35=5):} Graceful session termination.
\end{enumerate}
Integrity was maintained throughout; any tampering with the AES auth tag or the FIX checksum resulted in an immediate session reject.

\section{Conclusion}
The integration of TLS and AES-256-GCM provides a defense-in-depth strategy for financial messaging. While TLS secures the pipe, AES-GCM secures the data itself, ensuring that even in complex multi-hop routing environments, the financial instructions remain confidential and authentic.

\bibliographystyle{plain}
\bibliography{references}

\end{document}
